% ==============================================================
\section{提案手法}\label{sec:method}
% ==============================================================

本節では\textbf{イベント帰属パイプライン}の5ステップを詳述する．
Phase 1（Apriori-window）で得られたパターン集合 $\mathcal{F}$ と
その密集区間，および外部イベント集合 $\mathcal{E}$ を入力として，
統計的に有意なパターン--イベント帰属の集合を出力する．
アルゴリズム~\ref{alg:pipeline}に全体手順を示す．

\paragraph{前提条件: $L = 2$}
本パイプラインは2-itemset（$|P| = 2$）を前提として設計されている．
この制約は，(1) Step 5 のアイテム重複排除において，
2-itemset 間のアイテム共有グラフの連結成分分析が
副次パターンの同定に必要十分な構造を提供すること，
(2) Proposition~\ref{prop:dedup} の証明が $|P| = 2$ の場合に
$\beta^2$ vs $\beta \cdot p_{\text{base}}$ の比較に帰着されること，
に起因する．
$|P| \geq 3$ への拡張は，共有アイテム数に応じたスコア補正が必要となり，
今後の課題とする．

\begin{algorithm}[!ht]
\caption{帰属スコアリングと統計検定（Step 1--4）}
\label{alg:scoring}
\begin{algorithmic}[1]
\Require 共起パターン $\mathcal{F}\;(|P| \geq 2)$, 密集区間 $\{\mathcal{I}_P\}$, イベント $\mathcal{E}$, 設定 $(\theta, \sigma, B, \alpha, \delta)$
\Ensure 有意帰属集合 $\mathcal{R}$
\Statex \textbf{// Step 1--2: 変化点抽出 + フィルタリング}
\State $\text{RawResults} \gets []$
\For{各パターン $P \in \mathcal{F}$ with $\mathcal{I}_P \neq \emptyset$}
    \If{$\max(s_P) - \min(s_P) < \delta$} \textbf{continue} \EndIf \Comment{振幅フィルタ}
    \State $\mathcal{C}_P \gets \{s_i \text{(上昇)},\; e_i{+}1 \text{(下降)} \mid [s_i, e_i] \in \mathcal{I}_P\}$
    \State 各 $\tau \in \mathcal{C}_P$ の変化量 $\text{magnitude}(\tau)$ を計算 \Comment{式~(\ref{eq:magnitude})}
    \State $\mathcal{C}_P \gets \{\tau \in \mathcal{C}_P \mid \text{magnitude}(\tau) \geq \theta_{\text{mag}}\}$ \Comment{変化量フィルタ}
    \If{$\mathcal{C}_P = \emptyset$} \textbf{continue} \EndIf
    \Statex \textbf{// Step 3: 帰属スコアリング}
    \For{各イベント $e \in \mathcal{E}$}
        \State $A_{\text{obs}} \gets \sum_{\tau \in \mathcal{C}_P} A(P, \tau, e)$ \Comment{式~(\ref{eq:attribution_score})}
        \Statex \textbf{// Step 4: 円形シフト置換検定}
        \State $\text{count} \gets 0$
        \For{$b = 1$ to $B$}
            \State $\mathcal{E}_b \gets \Call{CircularShift}{\mathcal{E}, \delta_b}$
            \State $A_b \gets \sum_{\tau \in \mathcal{C}_P} A(P, \tau, e_b)$
            \If{$A_b \geq A_{\text{obs}}$} $\text{count} \gets \text{count} + 1$ \EndIf
        \EndFor
        \State $p_{\text{raw}} \gets (1 + \text{count}) / (B + 1)$ \Comment{式~(\ref{eq:pvalue})}
        \State $\text{RawResults}.\text{append}(P, e, A_{\text{obs}}, p_{\text{raw}})$
    \EndFor
\EndFor
\Statex \textbf{// Step 4 (続): グローバル BH 補正}
\State RawResults を $p_{\text{raw}}$ 昇順にソート
\State $M \gets |\text{RawResults}|$
\For{$i = M$ downto $1$}
    \State $\tilde{p}_{(i)} \gets \min(\tilde{p}_{(i+1)},\; \min(1,\; p_{(i)} \cdot M / i))$ \Comment{式~(\ref{eq:bh})}
\EndFor
\State $\mathcal{R} \gets \{(P, e) \mid \tilde{p}_{(i)} < \alpha\}$
\State \Return $\mathcal{R}$
\end{algorithmic}
\end{algorithm}


\subsection{Step 1: 密集区間からの変化点抽出}

Phase 1（Apriori-window）は各パターン $P \in \mathcal{F}$（$|P| \geq 2$）に対して，
サポートが閾値 $\theta$ 以上となるウィンドウ左端の連続区間（密集区間）
$\{[s_1, e_1], [s_2, e_2], \ldots\}$ を既に出力している．
本パイプラインではサポート時系列全体を再計算せず，これらの密集区間を直接利用する．
各密集区間 $[s_i, e_i]$ から上昇変化点 $s_i$ と下降変化点 $e_i + 1$（存在する場合）を抽出し，
変化量（式~(\ref{eq:magnitude})）の計算に必要な局所サポートのみを，
パターン出現リスト $T_P$ 上の二分探索により $O(\log |T_P|)$ で求める．

\subsection{Step 2: 変化量の評価}

Step 1 で密集区間の端点から変化点が既に抽出されているため，
本ステップではサポート時系列全体を走査する必要はない．
本手法が帰属対象とするのは
「密集区間の出現・消失に対応する切替点」であり，
密状態の内部で起こる任意のサポート変動ではない．

各変化点 $\tau$ に対して，変化量を交差前後の窓平均の差（レベルシフト量）として計算する:
\begin{equation}
\text{magnitude}(\tau) = |\bar{s}^{+}(\tau) - \bar{s}^{-}(\tau)|
\label{eq:magnitude}
\end{equation}
ここで $\bar{s}^{+}(\tau), \bar{s}^{-}(\tau)$ は定義~\ref{def:change_point} の局所平均であり，
$h$ はレベルウィンドウ幅（デフォルト $h = 20$）である．
この方式は瞬時差分（$s_P(\tau) - s_P(\tau-1)$）と比較して，
一過性の変動ではなく持続的な変化を捉える点で頑健性が高い．
すなわち，本手法では\textbf{検出位置}は閾値交差で決め，
\textbf{変化の強さ}は交差前後のレベル差で評価する．
なお，CUSUM のような一般的変化点検出器は比較対象として
\ref{sec:experiments}節で別途評価するが，主パイプラインには含めない．

\textbf{フィルタリング}: 有意性検定前の仮説空間削減のため，2つのフィルタを適用する:
\begin{itemize}
\item \emph{サポート振幅フィルタ}: $\max(s_P) - \min(s_P) < \delta$ のパターンを除外する．
ここで $\delta$ はユーザ指定の最小サポート振幅である．
このフィルタは時間的変動の乏しいパターンを排除し，仮説数を大幅に削減する．
\item \emph{変化量フィルタ}: 変化量が閾値未満の変化点を除外する．
\end{itemize}

\subsection{Step 3: 帰属スコアリング}

パターン $P$ の変化点 $\tau$ と各イベント $e \in \mathcal{E}$ に対して，
時間的関連を定量化する\emph{帰属スコア}を計算する:
\begin{equation}
A(P, \tau, e) = \text{prox}(\tau, e) \cdot \text{dir}(\tau, e) \cdot \text{mag}(\tau)
\label{eq:attribution_score}
\end{equation}

3つの構成要素は以下の通りである．

\textbf{近接度}: 指数カーネルによる時間的近さの指標:
\begin{equation}
\text{prox}(\tau, e) = \exp\left(-\frac{d(\tau, e)}{\sigma}\right)
\label{eq:proximity}
\end{equation}
ここで $d(\tau, e) = \min(|\tau - t_s|, |\tau - t_e|)$ は
変化点からイベント境界までの最小距離，$\sigma > 0$ は減衰パラメータである．

\textbf{方向整合性}: イベントの開始時にサポートが上昇し，
終了時に下降するという期待に基づく重み:
\begin{equation}
\text{dir}(\tau, e) = \begin{cases}
1.0 & \tau \geq t_s \text{ かつ方向が「上昇」} \\
1.0 & \tau \geq t_e \text{ かつ方向が「下降」} \\
0.5 & \text{その他（イベント前の変化等）} \\
\end{cases}
\label{eq:direction}
\end{equation}

\textbf{変化量}: 式~(\ref{eq:magnitude})のレベルシフト量．

パターン $P$ のイベント $e$ に対する観測スコアは，
全変化点にわたる帰属スコアの総和
$A_{\text{obs}}(P, e) = \sum_{\tau} A(P, \tau, e)$ とする．

\subsection{Step 4: 置換検定とグローバルFDR補正}

\subsubsection{円形シフト帰無分布}

パターン $P$ の変化点とイベント $e$ の時間的関連が統計的に有意かを検定するため，
円形シフト置換検定を用いる．
帰無仮説 $H_0$:「変化点の位置はイベント時刻と無関係である」のもと，
$B$ 回の置換を実行する．
各置換 $b$ では，ランダムなオフセット $\delta_b \sim \text{Uniform}\{1, \ldots, T-1\}$
だけ全イベントの時刻を一斉にずらす（円形シフト）:
\begin{equation}
e_b = (\textit{id}, \textit{name}, (t_s + \delta_b) \bmod T, (t_s + \delta_b + \Delta) \bmod T)
\end{equation}
ここで $\Delta = t_e - t_s$ はイベント持続時間，$T$ は時系列長である．

各置換に対してスコア $A_{\text{perm}}^{(b)}(P, e)$ を再計算し，
未補正 $p$ 値を以下で求める:
\begin{equation}
p_{\text{raw}}(P, e) = \frac{1 + |\{b : A_{\text{perm}}^{(b)}(P, e) \geq A_{\text{obs}}(P, e)\}|}{B + 1}
\label{eq:pvalue}
\end{equation}

\begin{proposition}[置換検定の正当性]\label{prop:perm_exact}
帰無仮説 $H_0$ のもとでは，観測スコア $A_{\text{obs}}$ と
$B$ 個の置換スコア $A_{\text{perm}}^{(1)}, \ldots, A_{\text{perm}}^{(B)}$ は交換可能（exchangeable）である．
したがって式~(\ref{eq:pvalue})の $p$ 値は
$\{1/(B+1), 2/(B+1), \ldots, 1\}$ 上の離散一様分布に従い，
任意の水準 $\alpha$ で第一種過誤率が $\alpha$ 以下に制御される~\cite{davison1997bootstrap}．
\end{proposition}

\subsubsection{グローバル Benjamini--Hochberg 補正}

$M = |\mathcal{F}'| \times |\mathcal{E}|$ 個の仮説を同時に検定するため，
全パターン・全イベントにわたり Benjamini--Hochberg (BH)
手続き~\cite{benjamini1995controlling}をグローバルに適用する．

未補正 $p$ 値を昇順にソートし $p_{(1)} \leq p_{(2)} \leq \cdots \leq p_{(M)}$ とする．
Step-down 調整済み $p$ 値を以下で計算する:
\begin{equation}
\tilde{p}_{(i)} = \min\!\left(\tilde{p}_{(i+1)},\; \min\!\left(1,\; p_{(i)} \cdot \frac{M}{i}\right)\right)
\label{eq:bh}
\end{equation}
ただし $\tilde{p}_{(M)} = \min(1, p_{(M)})$ とする．
$\tilde{p}_{(i)} < \alpha$ の仮説を棄却（帰属が有意）と判定する．

BH手続きは偽発見率（FDR）を水準 $\alpha$ 以下に保ちつつ，
Bonferroni補正と比べてより多くの真の帰属を検出できる（検出力が高い）．

\begin{remark}[PRDS 条件の充足]\label{rem:prds}
BH手続きが独立でない検定統計量に対してもFDRを制御するには，
$p$ 値が部分集合上の正回帰依存性（PRDS）条件を満たす必要がある~\cite{benjamini2001control}．
本設定ではアイテムを共有するパターン対のサポート時系列は正の相関を持つため，
対応する $p$ 値間に正の依存性が誘導され，PRDS条件は実務的に満たされると考えられる．
ただし，厳密な証明はデータ生成過程に依存するため，
本研究ではこの条件を実務的妥当性に基づいて適用する．
\end{remark}

\subsection{Step 5: アイテム重複排除}\label{sec:dedup}

本ステップは $|P|=2$ を前提とする．

\subsubsection{副次パターン問題}

外部イベントがアイテム $i$ の出現頻度を増加させると，
$i$ を含む\emph{全て}のパターンのサポートが相関して増加する．
例えば，イベントがアイテム $\{a, b\}$ をブーストし真のパターンが $(a, b)$ の場合，
$(c, a)$, $(d, b)$, $(e, a)$ などの副次パターンも
共有アイテム経由でサポートが連動変動するため有意性検定を通過する．

この問題は統計検定だけでは解消できない．
副次パターンも共有アイテム経由で実際にサポートが変動しているため，
置換検定でも有意と判定されてしまうからである．

\subsubsection{Union-Find クラスタリングによる重複排除}

この問題を以下の後処理ステップで解決する:

\begin{enumerate}
\item 各イベント $e$ について，$e$ に有意に帰属された全パターン $\{P_1, \ldots, P_n\}$ を収集する．
\item アイテム共有グラフを構築する: パターンをノード，少なくとも1つのアイテムを共有するパターン間にエッジを張る．
\item Union-Find（素集合）データ構造で連結成分を求める．
\item 各連結成分内で帰属スコア $A_{\text{obs}}(P, e)$ が最大のパターンのみを残す．
\end{enumerate}

\begin{proposition}[重複排除の正当性]\label{prop:dedup}
以下の簡約的な独立出現モデルを仮定する：
各アイテムはトランザクションごとに独立に出現し，
イベント期間外のベースライン出現確率は $p_{\text{base}}$ である．
真のパターン $P^* = \{i_1, i_2\}$ がイベント $e$ によりブースト確率 $\beta$ で
サポートを増加させられているとする．
副次パターン $P' = \{i_1, j\}$（$j \notin P^*$）において，
$j$ のベースライン出現確率を $p_{\text{base}}$ とする．
$\beta > p_{\text{base}}$ であれば，
$\text{mag}(\tau, P^*) > \text{mag}(\tau, P')$ が成立し，
$P^*$ と $P'$ はアイテム $i_1$ の共有により同一連結成分に属するため，
$P^*$ が保持され $P'$ は除去される．
\end{proposition}

\begin{proof}[証明の概略]
イベント $e$ の影響下では $i_1, i_2$ の各出現確率が $\beta$ だけ増加する．
$P^*$ のサポート増分は $\Delta s_{P^*} \propto \beta^2$（両アイテムがブースト）である．
一方 $P'$ のサポート増分は $\Delta s_{P'} \propto \beta \cdot p_{\text{base}}$
（$i_1$ はブーストされるが $j$ は独立にベースライン確率で出現）である．
$\beta > p_{\text{base}}$ のとき $\beta^2 > \beta \cdot p_{\text{base}}$ となり
$\text{mag}(\tau, P^*) > \text{mag}(\tau, P')$ が成立する．
実データにおいてブースト確率 $\beta$ が
個別アイテムのベースライン出現確率 $p_{\text{base}}$ を超える条件は，
イベントが検出可能な影響を持つ限り広く成立する．
\end{proof}

重複排除ステップの計算量はイベントあたり $O(n \cdot k)$ であり，
$n$ は有意帰属数，$k$ はパターン最大長である．
置換検定ステップに比して無視可能なコストである．

\subsection{計算量分析}

パイプライン全体の時間計算量は以下の通りである:
\begin{itemize}
\item \textbf{Step 1}（変化点抽出＋局所サポート計算）: $O(|\mathcal{F}| \cdot \bar{C} \cdot h \cdot \log n)$，$\bar{C}$ は平均密集区間数，$h$ はレベルウィンドウ幅，$n$ は平均パターン出現数
\item \textbf{Step 2}（変化量評価）: Step 1 に含まれる（追加走査なし）
\item \textbf{Step 3--4}（帰属＋置換検定）: $O(|\mathcal{F}'| \cdot |\mathcal{E}| \cdot B \cdot C)$，
$|\mathcal{F}'|$ はフィルタ後パターン数，$B$ は置換回数，$C$ は平均変化点数
\item \textbf{Step 5}（重複排除）: $O(R \cdot k)$，$R$ は有意結果数
\end{itemize}

サポート振幅フィルタが $|\mathcal{F}'|$ を90\%以上削減するため，
大規模パターン集合に対してもStep 3--4は効率的に動作する．
